% GENERATED FILE. Edit canonical lessons, then run npm run generate:books.
% canonical-source-hash: fnv1a64:f5c655879f0f1633
% canonical-lessons: KA-C187-sankhye, KA-C187-kamanabillu, KA-C187-hima, KA-C187-neralu, KA-C187-kattale

\chapter{Number, Rainbow, Snow, Shade, Darkness}
\label{ch:ka-number-rainbow-snow-shade-darkness}
\hlchaptermodality{187}

\begin{chapteropening}
\textbf{By the end of this chapter:} \emph{I can say a number, a rainbow, snow, shade and darkness.}

Complete the last lesson of chapter 187: I can say a number, a rainbow, snow, shade and darkness.
\end{chapteropening}

\section[\texorpdfstring{\kn{ಸಂಖ್ಯೆ} (saṅkhye)}{ಸಂಖ್ಯೆ (saṅkhye)}]{\kn{ಸಂಖ್ಯೆ} (saṅkhye) — a number}

\label{lesson:KA-C187-sankhye}



\begin{quote}
Before the new one: say the Kannada for a file, then the Kannada for a document.
\end{quote}

\subsection*{You'll want to know: \kn{ಸಂಖ್ಯೆ}}
\textbf{\kn{ಸಂಖ್ಯೆ}} — \emph{saṅkhye} — \textquotedblleft{}a number\textquotedblright{}.

\begin{grammarlens}[title={What you've built}]
One more word for this chapter.
\end{grammarlens}

\subsection*{Your turn}
\begin{itemize}
\raggedright
  \item \emph{Say it:} \emph{saṅkhye}
  \item \emph{Say it:} \emph{saṅkhye}, once more
  \item \emph{Say it:} say \emph{saṅkhye}
  \item \emph{Recall:} say the Kannada for a file, then the Kannada for a document, then say \emph{saṅkhye} again
\end{itemize}

\begin{tcolorbox}[breakable,colback=teal!4,colframe=teal!35!black,title={Before you move on}]
What does \kn{ಸಂಖ್ಯೆ} mean? (\textquotedblleft{}A number\textquotedblright{}.) Say it once more.
\end{tcolorbox}
\section[\texorpdfstring{\kn{ಕಾಮನಬಿಲ್ಲು} (kāmanabillu)}{ಕಾಮನಬಿಲ್ಲು (kāmanabillu)}]{\kn{ಕಾಮನಬಿಲ್ಲು} (kāmanabillu) — a rainbow}

\label{lesson:KA-C187-kamanabillu}



\begin{quote}
Before the new one: say the Kannada for a document, then the Kannada for a number.
\end{quote}

\subsection*{You'll want to know: \kn{ಕಾಮನಬಿಲ್ಲು}}
\textbf{\kn{ಕಾಮನಬಿಲ್ಲು}} — \emph{kāmanabillu} — \textquotedblleft{}a rainbow\textquotedblright{}.

\begin{grammarlens}[title={What you've built}]
One more word for this chapter.
\end{grammarlens}

\subsection*{Your turn}
\begin{itemize}
\raggedright
  \item \emph{Say it:} \emph{kāmanabillu}
  \item \emph{Say it:} \emph{kāmanabillu}, once more
  \item \emph{Say it:} say \emph{kāmanabillu}
  \item \emph{Recall:} say the Kannada for a document, then the Kannada for a number, then say \emph{kāmanabillu} again
\end{itemize}

\begin{tcolorbox}[breakable,colback=teal!4,colframe=teal!35!black,title={Before you move on}]
What does \kn{ಕಾಮನಬಿಲ್ಲು} mean? (\textquotedblleft{}A rainbow\textquotedblright{}.) Say it once more.
\end{tcolorbox}
\section[\texorpdfstring{\kn{ಹಿಮ} (hima)}{ಹಿಮ (hima)}]{\kn{ಹಿಮ} (hima) — snow}

\label{lesson:KA-C187-hima}



\begin{quote}
Before the new one: say the Kannada for a number, then the Kannada for a rainbow.
\end{quote}

\subsection*{You'll want to know: \kn{ಹಿಮ}}
\textbf{\kn{ಹಿಮ}} — \emph{hima} — \textquotedblleft{}snow\textquotedblright{}.

\begin{grammarlens}[title={What you've built}]
One more word for this chapter.
\end{grammarlens}

\subsection*{Your turn}
\begin{itemize}
\raggedright
  \item \emph{Say it:} \emph{hima}
  \item \emph{Say it:} \emph{hima}, once more
  \item \emph{Say it:} say \emph{hima}
  \item \emph{Recall:} say the Kannada for a number, then the Kannada for a rainbow, then say \emph{hima} again
\end{itemize}

\begin{tcolorbox}[breakable,colback=teal!4,colframe=teal!35!black,title={Before you move on}]
What does \kn{ಹಿಮ} mean? (\textquotedblleft{}Snow\textquotedblright{}.) Say it once more.
\end{tcolorbox}
\section[\texorpdfstring{\kn{ನೆರಳು} (neraḷu)}{ನೆರಳು (neraḷu)}]{\kn{ನೆರಳು} (neraḷu) — shade, a shadow}

\label{lesson:KA-C187-neralu}



\begin{quote}
Before the new one: say the Kannada for a rainbow, then the Kannada for snow.
\end{quote}

\subsection*{You'll want to know: \kn{ನೆರಳು}}
\textbf{\kn{ನೆರಳು}} — \emph{neraḷu} — \textquotedblleft{}shade, a shadow\textquotedblright{}.

\begin{grammarlens}[title={What you've built}]
One more word for this chapter.
\end{grammarlens}

\subsection*{Your turn}
\begin{itemize}
\raggedright
  \item \emph{Say it:} \emph{neraḷu}
  \item \emph{Say it:} \emph{neraḷu}, once more
  \item \emph{Say it:} say \emph{neraḷu}
  \item \emph{Recall:} say the Kannada for a rainbow, then the Kannada for snow, then say \emph{neraḷu} again
\end{itemize}

\begin{tcolorbox}[breakable,colback=teal!4,colframe=teal!35!black,title={Before you move on}]
What does \kn{ನೆರಳು} mean? (\textquotedblleft{}Shade, a shadow\textquotedblright{}.) Say it once more.
\end{tcolorbox}
\section[\texorpdfstring{\kn{ಕತ್ತಲೆ} (kattale)}{ಕತ್ತಲೆ (kattale)}]{\kn{ಕತ್ತಲೆ} (kattale) — darkness}

\label{lesson:KA-C187-kattale}



\begin{quote}
Before the new one: say the Kannada for snow, then the Kannada for shade.
\end{quote}

\subsection*{You'll want to know: \kn{ಕತ್ತಲೆ}}
\textbf{\kn{ಕತ್ತಲೆ}} — \emph{kattale} — \textquotedblleft{}darkness\textquotedblright{}.

\begin{grammarlens}[title={What you've built}]
One more word for this chapter.
\end{grammarlens}

\subsection*{Your turn}
\begin{itemize}
\raggedright
  \item \emph{Say it:} \emph{kattale}
  \item \emph{Say it:} \emph{kattale}, once more
  \item \emph{Say it:} say \emph{kattale}
  \item \emph{Recall:} say the Kannada for snow, then the Kannada for shade, then say \emph{kattale} again
\end{itemize}

\begin{tcolorbox}[breakable,colback=teal!4,colframe=teal!35!black,title={Before you move on}]
What does \kn{ಕತ್ತಲೆ} mean? (\textquotedblleft{}Darkness\textquotedblright{}.) Say it once more.
\end{tcolorbox}
